\section{Quản lý Xuất Nhập File (File IO)}

\subsection{Tạo, Xem và Thêm nội dung vào File}

\textbf{Cách thực hiện:}
Sử dụng toán tử redirect (\texttt{>}) để tạo file \texttt{myfile.txt} chứa nội dung "Hello, Linux!" và xem nội dung bằng \texttt{cat}. Tiếp tục dùng toán tử append (\texttt{>>}) để thêm dòng mới mà không làm mất dữ liệu cũ.

\begin{lstlisting}[language=bash]
# Tao file va xem
echo "Hello, Linux!" > myfile.txt
cat myfile.txt

# Them noi dung va xem lai
echo "Welcome to Linux file management." >> myfile.txt
cat myfile.txt
\end{lstlisting}

\textbf{Kết quả:}
Tập tin được tạo và thêm nội dung thành công, các toán tử redirect hoạt động đúng như lý thuyết.

\begin{figure}[H]
    \centering
    \begin{subfigure}[b]{0.45\textwidth}
        \centering
        \includegraphics[width=\textwidth]{03/create-myfile.png}
        \caption{Tạo file myfile.txt}
    \end{subfigure}
    \hfill
    \begin{subfigure}[b]{0.45\textwidth}
        \centering
        \includegraphics[width=\textwidth]{03/append-myfile.png}
        \caption{Thêm nội dung (Append)}
    \end{subfigure}
    \caption{Thao tác cơ bản với File}
\end{figure}

\subsection{Đếm số dòng và từ}

\textbf{Cách thực hiện:}
Sử dụng lệnh \texttt{wc} (word count) với tùy chọn \texttt{-l} để đếm số dòng và \texttt{-w} để đếm số từ trong tập tin \texttt{myfile.txt}.

\begin{lstlisting}[language=bash]
wc -l myfile.txt
wc -w myfile.txt
\end{lstlisting}

\textbf{Kết quả:}
Hệ thống đếm chính xác số lượng dòng (2 dòng) và số lượng từ.

\begin{figure}[H]
    \centering
    \begin{subfigure}[b]{0.45\textwidth}
        \centering
        \includegraphics[width=\textwidth]{03/wc-lines.png}
        \caption{Đếm số dòng}
    \end{subfigure}
    \hfill
    \begin{subfigure}[b]{0.45\textwidth}
        \centering
        \includegraphics[width=\textwidth]{03/wc-words.png}
        \caption{Đếm số từ}
    \end{subfigure}
    \caption{Lệnh đếm dòng và từ}
\end{figure}

\subsection{Sao chép, Di chuyển và Xóa File}

\textbf{Cách thực hiện:}
Tiến hành sao chép \texttt{myfile.txt} thành \texttt{myfile\_backup.txt}, sau đó di chuyển bản backup vào thư mục \texttt{/tmp/}. Cuối cùng, xóa bỏ file gốc \texttt{myfile.txt}.

\begin{lstlisting}[language=bash]
# Sao chep
cp myfile.txt myfile_backup.txt

# Di chuyen
mv myfile_backup.txt /tmp/

# Xoa file
rm myfile.txt
\end{lstlisting}

\textbf{Kết quả:}
Các lệnh thao tác tập tin hoạt động chính xác. File gốc được xóa và bản backup nằm an toàn trong \texttt{/tmp/}.

\begin{figure}[H]
    \centering
    \begin{subfigure}[b]{0.45\textwidth}
        \centering
        \includegraphics[width=\textwidth]{03/cp-myfile-backup.png}
        \caption{Sao chép file}
    \end{subfigure}
    \hfill
    \begin{subfigure}[b]{0.45\textwidth}
        \centering
        \includegraphics[width=\textwidth]{03/mv-to-tmp.png}
        \caption{Di chuyển file sang /tmp}
    \end{subfigure}
    
    \vspace{0.5cm}
    \begin{subfigure}[b]{0.45\textwidth}
        \centering
        \includegraphics[width=\textwidth]{03/rm-myfile.png}
        \caption{Xóa file gốc}
    \end{subfigure}
    \caption{Thao tác quản lý File nâng cao}
\end{figure}

\subsection{Tìm kiếm, Đổi tên và Xem từng phần}

\textbf{Cách thực hiện:}
- Dùng \texttt{grep} để tìm kiếm chuỗi "Linux" trong file backup.
- Dùng \texttt{mv} để đổi tên file thành \texttt{linuxfile.txt}.
- Tạo thư mục \texttt{backup} và di chuyển file vào trong.
- Dùng \texttt{less} để xem nội dung file.

\begin{lstlisting}[language=bash]
# Tim kiem
grep "Linux" /tmp/myfile_backup.txt

# Doi ten va tao thu muc
mv /tmp/myfile_backup.txt /tmp/linuxfile.txt
mkdir backup
mv /tmp/linuxfile.txt backup/

# Xem noi dung tung trang
less backup/linuxfile.txt
\end{lstlisting}

\textbf{Kết quả:}
Công cụ \texttt{grep} tìm được đúng dòng chứa từ khóa. \texttt{mv} và \texttt{mkdir} tạo cấu trúc thư mục thành công.

\begin{figure}[H]
    \centering
    \begin{subfigure}[b]{0.45\textwidth}
        \centering
        \includegraphics[width=\textwidth]{03/grep-linux.png}
        \caption{Tìm kiếm với grep}
    \end{subfigure}
    \hfill
    \begin{subfigure}[b]{0.45\textwidth}
        \centering
        \includegraphics[width=\textwidth]{03/rename-linuxfile.png}
        \caption{Đổi tên file}
    \end{subfigure}
    
    \vspace{0.5cm}
    \begin{subfigure}[b]{0.45\textwidth}
        \centering
        \includegraphics[width=\textwidth]{03/mkdir-backup-mv.png}
        \caption{Tạo thư mục và di chuyển}
    \end{subfigure}
    \hfill
    \begin{subfigure}[b]{0.45\textwidth}
        \centering
        \includegraphics[width=\textwidth]{03/less-linuxfile.png}
        \caption{Xem file bằng less}
    \end{subfigure}
    \caption{Các công cụ tìm kiếm và điều hướng}
\end{figure}
